\section{实验过程与结果分析}\label{sec:experiments}
\subsection{实验安排与检查结果}
我们先进行程序和单机调试，再逐步连接相机、底盘及多车任务。表中汇总了各部分完成的检查和仍需补充的试验。程序能够运行只是第一步，实际效果还需要结合现场画面和重复测试来评价。

\begin{longtable}{p{27mm}p{58mm}p{53mm}}
\caption{各部分实验检查情况}\\\toprule
项目 & 检查内容与结果 & 进一步测试\\\midrule\endfirsthead
\toprule 项目 & 检查内容与结果 & 进一步测试\\\midrule\endhead
动作组 & 核对组合动作的帧序、时长和手臂位置 & 满载步态、抓取及放置稳定性\\
标记路线 & 完成标记识别与局部路线调试，走向A并开始左转 & A/B/C完整路线重复运行\\
相机输入 & 重新接插后恢复取流，连续取得640$\times$480画面 & 舵机负载变化时的断线情况\\
基础底盘 & 启动与键盘遥控可用 & 相机、扫描和导航联合运行\\
控制逻辑 & 11项离线控制与停止检查通过 & 实际通信、底盘响应和停止延迟\\
三车队列 & 完成直行、圆弧转弯和顺序停车 & 专门障碍、丢失与重捕获场景\\
大组联调 & 结合现场视频检查前后车及转角现象 & 各方式同轮运行与完整终点记录\\\bottomrule
\end{longtable}

\subsection{分阶段实验与异常注入}
\begin{enumerate}
\item \textbf{单动作与视觉。}人工监督下分别测试站立、抓取、抱物前进、左右转和放置；固定相机姿态后测量标记距离/方位，统计多次误差与失效距离。
\item \textbf{人形单段与完整路线。}先测直行偏移，再测A左转/B右转，最后无人工干预完成A/B/C全程。分别记录动作次数、转角偏差、接触、掉落及相机断线。
\item \textbf{地图和单车导航。}确认35\,cm障碍进入原始深度与扫描，再检查地图。先静态目标导航、定位丢失和规划失败，再接入移动AR目标。
\item \textbf{目标与任务失效。}遮挡标签、停止目标消息、停止任务心跳、断开扫描或使TF过期，逐项测量停止延迟。人形最新程序的盲行与近距丢码启发式应单独评估。
\item \textbf{多车联调。}固定队列拓扑、共享地图与时间基准，避免速度控制源冲突；记录各车到达终点条件，不只记录领行机器人结束。
\end{enumerate}

\subsection{指标定义}
距离标定用均方根误差$\operatorname{RMSE}_d=\sqrt{N^{-1}\sum_i(\hat d_i-d_i)^2}$；$d_i$必须是与PnP输出同一定义的真实距离。跟随间距误差定义为$e_f(t)=d(t)-d_0$，报告平均绝对误差与最大偏差，同时说明$d$采用欧氏距离还是路径弧长。完整任务成功率定义为$N_{\rm success}/N_{\rm trials}$，每次试验均需记录，不能剔除失败试跑而不说明。

对接触次数、掉落次数、目标丢失持续时间、最长无新帧时间、停止延迟和各车到达时间分别统计。转角是否达到路线要求应通过外部测量或可复核视频判定，不能把“看见下一张码”直接记为精确90度。

\begin{table}[H]\centering\small
\caption{后续量化试验安排}
\begin{tabularx}{\linewidth}{p{42mm}YY}\toprule
项目 & 测量方法 & 用途\\\midrule
标记距离与方位 & 固定摆位、多距离和多角度比较 & 检查视觉判断与实际距离的关系\\
满载步态与转向 & 重复执行单步及左右转动作 & 统计步长、偏航与转角分布\\
完整路线 & 连续运行并记录失败轮次 & 比较路线完成率与运行时间\\
35\,cm障碍 & 同时查看深度、扫描及实际运动 & 排查漏检与避让效果\\
故障停车 & 分别中断目标、任务和扫描输入 & 测量发现异常到实际停稳的时间\\
全队终点 & 连续拍摄各车进入并停稳 & 检查队列到达与停车顺序\\\bottomrule
\end{tabularx}\end{table}

目前的记录以调试现象为主，量化比较还不充分。后续按表中方法重复试验，记录运行参数、日志和视频，并同时保留成功与失败的情况，才能判断一次调整是否有效。


\subsection{数据归档与照片使用}
运行记录模板保存在\code{latex/data/experiment_records.csv}。每行对应一次试验，保存代码版本、设备、运行参数、结果及证据路径；原始视频和bag较大时放在团队共享存储中，报告记录可访问的证据位置。不要将不同机器人的动作参数或不同版本的相机设置合并统计。

我们拍摄了通道、机器人和调试窗口，并从比赛视频中选取了运行画面，用来说明场地和程序判断之间的关系。下一节按视频时间介绍这些现象；单帧画面只反映某个时刻，完整过程还需要结合连续视频。
